\subsection{关于全纯函数演算的补充注记}

对 C 的每个子集 X，用 \(X^{*}\) 表示 X 在复共轭下的像。设 U 是 C 的一个开子集，\(g: U \to C\) 是一个全纯函数。于是函数 \(f^{*}: z \mapsto \overline{g(\overline{z})}\) 在 \(U^{*}\) 上有定义且全纯。映射 \(f \mapsto f^{*}\) 是从 \(\mathcal{O}(U)\) 到 \(\mathcal{O}(U^{*})\) 上的一个连续双射，使得 \((f_{1}f_{2})^{*} = f_{1}^{*}f_{2}^{*}\) 与 \((f_{1} + f_{2})^{*} = f_{1}^{*} + f_{2}^{*}\) 对 \(f_{1}\) 与 \(f_{2}\)（\(\mathcal{O}(U)\) 中的）成立。

设 C 是 C 的一个紧子集。映射 \(f \mapsto f^{*}\)（\(\mathcal{O}(U)\) 到 \(\mathcal{O}(U^{*})\) 中的）当 U 取遍 C 的包含 C 的开子集时，诱导出空间 \(\mathcal{O}(C)\) 到 \(\mathcal{O}(C^{*})\) 上的一个连续双射（I, 第 49 页, 第 1 目），它也记作 \(f \mapsto f^{*}\)，且满足 \((f_{1}f_{2})^{*} = f_{1}^{*}f_{2}^{*}\) 与 \((f_{1} + f_{2})^{*} = f_{1}^{*} + f_{2}^{*}\) 对 \(f_{1}\) 与 \(f_{2}\)（\(\mathcal{O}(C)\) 中的）。

命题 3. --- 设 A 是一个含单位元的对合巴拿赫代数。设 \(a \in \mathrm{A}\)。\(a^{*}\) 的谱是 \(\mathrm{Sp}_{\mathrm{A}}(a)^{*}\)，即 \(\mathrm{Sp}_{\mathrm{A}}(a)\) 在复共轭下的像。对每个 \(f \in \mathcal{O}(\mathrm{Sp}_{\mathrm{A}}(a))\)，我们有 \(f(a)^{*} = f^{*}(a^{*})\)。

第一个断言由 I 第 97 页得出。由前所述，映射 \(\varphi\)（\(\mathcal{O}(\mathrm{Sp}_{\mathrm{A}}(a))\) 到 A 中的）由 \(f \mapsto (f^{*}(a^{*}))^{*}\) 定义，它是 \(\mathcal{O}(\mathrm{Sp}_{\mathrm{A}}(a))\) 到 A 中的一个连续幺态射，使得在 \(\mathrm{Sp}_{\mathrm{A}}(a)\) 的一个邻域中 \(\mathbf{C}\) 的恒同函数的芽的像等于 \(a\)。因此，\(\varphi\) 就是全纯函数演算的映射 \(f \mapsto f(a)\)（I, 第 74 页, 定理 5）。

引理 1. --- 设 D 是 C 中的单位开圆盘。存在唯一的在 D 上定义的全纯函数 f，使得 \(f(z)^{2} = 1 - z\) 对所有 \(z \in D\)，且 \(f(0) = 1\)。我们有 \(f^{*} = f\)。

幂级数

\[
\sum_ {n = 0} ^ {+ \infty} \binom {1 / 2} {n} (- z) ^ {n}
\]

的收敛半径是 1，其和 \(f\) 定义了 D 上的一个全纯函数（VAR, R1, 第 27 页, 3.2.9），它在 0 处取值 1。它满足 \(f(x)=\sqrt{1-x}\) 对所有 \(x \in D \cap R\)（FVR, III, 第 19 页），因而 \(f(z)^{2} = 1 - z\) 对 \(z \in D\)，因为差 \(f(z)^{2} - (1 - z)\) 是一个各阶导数都在 0 处为零的全纯函数（VAR, R1, 第 27 页, 3.2.5）。按定义即可验证 \(f^{*} = f\)。

设 \(g\) 是 D 上的一个全纯函数，使得 \(g(z)^2 = 1 - z\) 对所有 \(z \in \mathrm{D}\)，且 \(g(0) = 1\)。函数 \(g\) 不取零值，于是 D 上的连续函数 \(f / g\) 取值于 \(\{-1, 1\}\)；由于 D 是连通的且 \(f(0) = g(0)\)，我们有 \(f = g\)。
% semantic-fragments: legacy-input-seam
\relax{}
